\scanpage{16}
% archive foliation: 15
\unsure{\deleted{Si $X$ est une courbe alg. de genre $g$,}}
\inserted{\unsure{Si $X$ courbe algébrique lisse} \inserted{propre} de genre $g$\deleted{$\unsure{\ge0}$}}, et $k$ de car. nulle, conditions équivalentes
\[
\begin{array}{ll}
d)&\dim H^1_{DR}(X,\unsure{\mathfrak g})=6g-6\quad\deleted{\unsure{(si $g$}}\\
d')&H^0_{DR}\deleted{\unsure{\xi}}(X,\unsure{\mathfrak g})\ (=H^{0\unsure{\xi|\zeta}}(X,\unsure{\mathfrak g}))=0\\
d'')&H^2_{DR}(X,\unsure{\mathfrak g})\ (=H^{2\unsure{\xi|\zeta}}(X,\unsure{\mathfrak g}))=0
\end{array}
\]
\margin{(conditions jamais vérifiées si $g=0$ -- et pour cause !)}

[\uline{NB} On a $\unsure{\mathfrak g}\simeq\check{\unsure{\mathfrak g}}$ (Killing) donc
\[
(*)\quad H^0_{DR}(X,\unsure{\mathfrak g})\simeq H^2_{DR}(X,\unsure{\mathfrak g})^\vee
\]
(\unsure{d'où équiv. de $d'$ et $d''$}) et formule de EP
\[
\begin{aligned}
(**)\quad\chi^{\unsure{\xi|\zeta}}_{DR}(X,\unsure{\mathfrak g})
&=\sum(-1)^i\dim H^i_{DR}(X,\unsure{\mathfrak g})
=\underbrace{3}_{\operatorname{rg}\unsure{\mathfrak g}}\ \underbrace{\chi^{\unsure{\xi|\zeta}}(X,\mathcal O_X)}_{2-2g}\\
&=\deleted{\unsure{6g}}6-6g
\end{aligned}
\]
donc, combinant avec $(*)$
\[
(***)\quad \dim H^1_{DR}(X,\unsure{\mathfrak g})=2\dim H^0_{DR}(X,\unsure{\mathfrak g})+(6g-6)
\]
\margin{\deleted{\unsure{dim gr. $G$}}}
d'où l'équivalence de $d')$ avec $d)$.

Enfin, on sait
\[
H^0_{DR}(X,\unsure{\mathfrak g})\simeq\uline{\operatorname{Lie}}\operatorname{Aut}^{\unsure{\xi|\zeta}}_{X/k}(P)
\]
d'où l'équivalence de $d')$ avec $b)$, et la formule intéressante
\[
\begin{aligned}
\dim\underbrace{H^1_{DR}(X,\unsure{\mathfrak g})}_{\substack{\text{espace tangent}\\\text{à l'espace modulaire}\\M_g^{\unsure{\xi|\zeta}}\text{ au pt correspondant}\\\text{à }(X,P,\unsure{\xi|\zeta})}}
&=2\dim\operatorname{Aut}^{\unsure{\xi|\zeta}}_{X/k}(P)+6g-6\\
&=2\bigl(\dim\operatorname{Aut}^{\unsure{\xi|\zeta}}_{X/k}(P)+3g-3\bigr)
\end{aligned}
\]

\uline{NB} La condition $d')$ implique que $M_g^{\unsure{\xi|\zeta}}$ ``n'a pas d'autom. infinitésimaux'', $a)$ qu'il ``n'a pas d'autom. finis'', donc ensemble que $M_g^{\unsure{\xi|\zeta}}$ est ``\uline{rigide}'' -- tout comme $M_g$ -- donc un foncteur représentable ! La condition $d'')$ implique que $M_g^{\unsure{\xi|\zeta}}$ est \inserted{form.} lisse sur $M_g$, et $d)$ que sa dim relative sur $M_g$ est $6g-6$.

Sans doute peut-on prouver la représentabilité relative (du moins au sens espaces alg. d'Artin)
